% =====================================================================
%  Star.org 白皮书 · 共享排版样式
% =====================================================================
%  用法：语言源文件先 % =====================================================================
%  Star.org 白皮书 · 共享排版样式
% =====================================================================
%  用法：语言源文件先 % =====================================================================
%  Star.org 白皮书 · 共享排版样式
% =====================================================================
%  用法：语言源文件先 % =====================================================================
%  Star.org 白皮书 · 共享排版样式
% =====================================================================
%  用法：语言源文件先 \input{preamble.tex}，再写正文。
%  字体使用仓库内置的 Noto Sans/Serif SC 子集（fonts/），
%  渲染不依赖运行环境是否安装中文字体（见 DECISIONS D10）。
% =====================================================================

\usepackage{fontspec}
\usepackage{xeCJK}

% ---- 字体：仓库内置的 Noto Sans/Serif SC 子集 ----
% 路径相对于编译工作目录（docs/whitepaper/）。字体随仓库内置，
% 渲染结果不再依赖运行环境是否安装中文字体（DECISIONS D10）。
% 子集覆盖 CJK 统一表意文字 U+4E00–U+9FFF 全区间（见 fonts/manifest.json）。
% 内置子集只有 Regular 一个字重，AutoFakeBold 用于合成粗体，
% 否则 \bfseries 会静默回退到 Regular、标题看不出层级。
\setCJKmainfont[Path=../../fonts/,AutoFakeBold=2.5]{NotoSerifSC-Regular.subset.otf}
\setCJKsansfont[Path=../../fonts/,AutoFakeBold=2.5]{NotoSansSC-Regular.subset.otf}
\setCJKmonofont[Path=../../fonts/]{NotoSansSC-Regular.subset.otf}
% 拉丁字母与数字沿用 TeX 自带的 Latin Modern（任何 TeX Live 环境都有），
% 不依赖系统是否安装 Helvetica 一类的西文字体。

\usepackage{geometry}
\usepackage{xcolor}
\usepackage{booktabs}
\usepackage{array}
\usepackage{tabularx}
\usepackage{longtable}
\usepackage{fancyhdr}
\usepackage{setspace}
\usepackage{parskip}
\usepackage{caption}
\usepackage{ragged2e}
\usepackage{microtype}
\usepackage[hidelinks]{hyperref}

% ---- 页面 ----
\geometry{a4paper, top=26mm, bottom=24mm, left=24mm, right=24mm}

% ---- 配色（与站点一致的深蓝 + 金）----
\definecolor{ink}{HTML}{111827}
\definecolor{inksoft}{HTML}{4B5563}
\definecolor{inkfaint}{HTML}{6B7280}
\definecolor{accent}{HTML}{1E3A8A}
\definecolor{gold}{HTML}{B08430}
\definecolor{line}{HTML}{D8DEE9}
\definecolor{surface}{HTML}{F6F7FB}

% ---- 页眉页脚 ----
\pagestyle{fancy}
\fancyhf{}
\renewcommand{\headrulewidth}{0.4pt}
\renewcommand{\footrulewidth}{0pt}
\fancyhead[L]{\footnotesize\color{inkfaint}\sffamily Star.org}
\fancyhead[R]{\footnotesize\color{inkfaint}\sffamily\WHTITLE}
\fancyfoot[C]{\footnotesize\color{inkfaint}\thepage}

% ---- 标题样式（不依赖 titlesec / sectsty）----
\makeatletter
\renewcommand{\section}{\@startsection{section}{1}{0pt}%
  {22pt plus 4pt minus 2pt}{10pt plus 2pt minus 2pt}%
  {\normalfont\sffamily\Large\bfseries\color{accent}}}
\renewcommand{\subsection}{\@startsection{subsection}{2}{0pt}%
  {16pt plus 3pt minus 2pt}{7pt plus 2pt minus 2pt}%
  {\normalfont\sffamily\large\bfseries\color{ink}}}
\renewcommand{\subsubsection}{\@startsection{subsubsection}{3}{0pt}%
  {12pt plus 2pt minus 2pt}{6pt plus 1pt minus 1pt}%
  {\normalfont\sffamily\normalsize\bfseries\color{inksoft}}}
\makeatother

% 章标题前不加首行缩进，段落间距由 parskip 控制
\setlength{\parindent}{0pt}

% ---- 引用块（核心主张）：左侧金色竖线 + 浅底 ----
\newcommand{\claim}[1]{%
  \par\vspace{6pt}%
  \noindent\begin{minipage}{\linewidth}%
    \setlength{\fboxsep}{12pt}%
    \colorbox{surface}{%
      \begin{minipage}{\dimexpr\linewidth-2\fboxsep\relax}%
        \color{accent}\bfseries\sffamily #1%
      \end{minipage}}%
  \end{minipage}\par\vspace{8pt}%
}

% ---- 一句话结语块 ----
\newcommand{\takeaway}[1]{%
  \par\vspace{8pt}\noindent\begin{minipage}{\linewidth}%
    \setlength{\fboxsep}{14pt}%
    \colorbox{accent}{%
      \begin{minipage}{\dimexpr\linewidth-2\fboxsep\relax}%
        \color{white}#1%
      \end{minipage}}%
  \end{minipage}\par\vspace{10pt}%
}

% ---- 表格统一样式 ----
\newcommand{\wtable}[1]{%
  \par\vspace{4pt}\noindent\small #1\par\vspace{6pt}%
}

% 目录条目字号
\renewcommand{\contentsname}{\sffamily\Large\bfseries\color{accent}\WHTOCNAME}

% 英文日期：ctexart 下 \today 输出中文格式（2026 年 9 月 21 日），英文版单独用这个
\newcommand{\WHdateEN}{%
  \ifcase\month\or January\or February\or March\or April\or May\or June\or
  July\or August\or September\or October\or November\or December\fi
  \space\the\day,\space\the\year}

% hyperref 书签与元数据
\hypersetup{
  colorlinks=true,
  linkcolor=accent,
  urlcolor=accent,
  pdftitle={\WHTITLE},
  pdfauthor={Star.org},
  pdfsubject={\WHSUBJECT},
}
，再写正文。
%  字体使用仓库内置的 Noto Sans/Serif SC 子集（fonts/），
%  渲染不依赖运行环境是否安装中文字体（见 DECISIONS D10）。
% =====================================================================

\usepackage{fontspec}
\usepackage{xeCJK}

% ---- 字体：仓库内置的 Noto Sans/Serif SC 子集 ----
% 路径相对于编译工作目录（docs/whitepaper/）。字体随仓库内置，
% 渲染结果不再依赖运行环境是否安装中文字体（DECISIONS D10）。
% 子集覆盖 CJK 统一表意文字 U+4E00–U+9FFF 全区间（见 fonts/manifest.json）。
% 内置子集只有 Regular 一个字重，AutoFakeBold 用于合成粗体，
% 否则 \bfseries 会静默回退到 Regular、标题看不出层级。
\setCJKmainfont[Path=../../fonts/,AutoFakeBold=2.5]{NotoSerifSC-Regular.subset.otf}
\setCJKsansfont[Path=../../fonts/,AutoFakeBold=2.5]{NotoSansSC-Regular.subset.otf}
\setCJKmonofont[Path=../../fonts/]{NotoSansSC-Regular.subset.otf}
% 拉丁字母与数字沿用 TeX 自带的 Latin Modern（任何 TeX Live 环境都有），
% 不依赖系统是否安装 Helvetica 一类的西文字体。

\usepackage{geometry}
\usepackage{xcolor}
\usepackage{booktabs}
\usepackage{array}
\usepackage{tabularx}
\usepackage{longtable}
\usepackage{fancyhdr}
\usepackage{setspace}
\usepackage{parskip}
\usepackage{caption}
\usepackage{ragged2e}
\usepackage{microtype}
\usepackage[hidelinks]{hyperref}

% ---- 页面 ----
\geometry{a4paper, top=26mm, bottom=24mm, left=24mm, right=24mm}

% ---- 配色（与站点一致的深蓝 + 金）----
\definecolor{ink}{HTML}{111827}
\definecolor{inksoft}{HTML}{4B5563}
\definecolor{inkfaint}{HTML}{6B7280}
\definecolor{accent}{HTML}{1E3A8A}
\definecolor{gold}{HTML}{B08430}
\definecolor{line}{HTML}{D8DEE9}
\definecolor{surface}{HTML}{F6F7FB}

% ---- 页眉页脚 ----
\pagestyle{fancy}
\fancyhf{}
\renewcommand{\headrulewidth}{0.4pt}
\renewcommand{\footrulewidth}{0pt}
\fancyhead[L]{\footnotesize\color{inkfaint}\sffamily Star.org}
\fancyhead[R]{\footnotesize\color{inkfaint}\sffamily\WHTITLE}
\fancyfoot[C]{\footnotesize\color{inkfaint}\thepage}

% ---- 标题样式（不依赖 titlesec / sectsty）----
\makeatletter
\renewcommand{\section}{\@startsection{section}{1}{0pt}%
  {22pt plus 4pt minus 2pt}{10pt plus 2pt minus 2pt}%
  {\normalfont\sffamily\Large\bfseries\color{accent}}}
\renewcommand{\subsection}{\@startsection{subsection}{2}{0pt}%
  {16pt plus 3pt minus 2pt}{7pt plus 2pt minus 2pt}%
  {\normalfont\sffamily\large\bfseries\color{ink}}}
\renewcommand{\subsubsection}{\@startsection{subsubsection}{3}{0pt}%
  {12pt plus 2pt minus 2pt}{6pt plus 1pt minus 1pt}%
  {\normalfont\sffamily\normalsize\bfseries\color{inksoft}}}
\makeatother

% 章标题前不加首行缩进，段落间距由 parskip 控制
\setlength{\parindent}{0pt}

% ---- 引用块（核心主张）：左侧金色竖线 + 浅底 ----
\newcommand{\claim}[1]{%
  \par\vspace{6pt}%
  \noindent\begin{minipage}{\linewidth}%
    \setlength{\fboxsep}{12pt}%
    \colorbox{surface}{%
      \begin{minipage}{\dimexpr\linewidth-2\fboxsep\relax}%
        \color{accent}\bfseries\sffamily #1%
      \end{minipage}}%
  \end{minipage}\par\vspace{8pt}%
}

% ---- 一句话结语块 ----
\newcommand{\takeaway}[1]{%
  \par\vspace{8pt}\noindent\begin{minipage}{\linewidth}%
    \setlength{\fboxsep}{14pt}%
    \colorbox{accent}{%
      \begin{minipage}{\dimexpr\linewidth-2\fboxsep\relax}%
        \color{white}#1%
      \end{minipage}}%
  \end{minipage}\par\vspace{10pt}%
}

% ---- 表格统一样式 ----
\newcommand{\wtable}[1]{%
  \par\vspace{4pt}\noindent\small #1\par\vspace{6pt}%
}

% 目录条目字号
\renewcommand{\contentsname}{\sffamily\Large\bfseries\color{accent}\WHTOCNAME}

% 英文日期：ctexart 下 \today 输出中文格式（2026 年 9 月 21 日），英文版单独用这个
\newcommand{\WHdateEN}{%
  \ifcase\month\or January\or February\or March\or April\or May\or June\or
  July\or August\or September\or October\or November\or December\fi
  \space\the\day,\space\the\year}

% hyperref 书签与元数据
\hypersetup{
  colorlinks=true,
  linkcolor=accent,
  urlcolor=accent,
  pdftitle={\WHTITLE},
  pdfauthor={Star.org},
  pdfsubject={\WHSUBJECT},
}
，再写正文。
%  字体使用仓库内置的 Noto Sans/Serif SC 子集（fonts/），
%  渲染不依赖运行环境是否安装中文字体（见 DECISIONS D10）。
% =====================================================================

\usepackage{fontspec}
\usepackage{xeCJK}

% ---- 字体：仓库内置的 Noto Sans/Serif SC 子集 ----
% 路径相对于编译工作目录（docs/whitepaper/）。字体随仓库内置，
% 渲染结果不再依赖运行环境是否安装中文字体（DECISIONS D10）。
% 子集覆盖 CJK 统一表意文字 U+4E00–U+9FFF 全区间（见 fonts/manifest.json）。
% 内置子集只有 Regular 一个字重，AutoFakeBold 用于合成粗体，
% 否则 \bfseries 会静默回退到 Regular、标题看不出层级。
\setCJKmainfont[Path=../../fonts/,AutoFakeBold=2.5]{NotoSerifSC-Regular.subset.otf}
\setCJKsansfont[Path=../../fonts/,AutoFakeBold=2.5]{NotoSansSC-Regular.subset.otf}
\setCJKmonofont[Path=../../fonts/]{NotoSansSC-Regular.subset.otf}
% 拉丁字母与数字沿用 TeX 自带的 Latin Modern（任何 TeX Live 环境都有），
% 不依赖系统是否安装 Helvetica 一类的西文字体。

\usepackage{geometry}
\usepackage{xcolor}
\usepackage{booktabs}
\usepackage{array}
\usepackage{tabularx}
\usepackage{longtable}
\usepackage{fancyhdr}
\usepackage{setspace}
\usepackage{parskip}
\usepackage{caption}
\usepackage{ragged2e}
\usepackage{microtype}
\usepackage[hidelinks]{hyperref}

% ---- 页面 ----
\geometry{a4paper, top=26mm, bottom=24mm, left=24mm, right=24mm}

% ---- 配色（与站点一致的深蓝 + 金）----
\definecolor{ink}{HTML}{111827}
\definecolor{inksoft}{HTML}{4B5563}
\definecolor{inkfaint}{HTML}{6B7280}
\definecolor{accent}{HTML}{1E3A8A}
\definecolor{gold}{HTML}{B08430}
\definecolor{line}{HTML}{D8DEE9}
\definecolor{surface}{HTML}{F6F7FB}

% ---- 页眉页脚 ----
\pagestyle{fancy}
\fancyhf{}
\renewcommand{\headrulewidth}{0.4pt}
\renewcommand{\footrulewidth}{0pt}
\fancyhead[L]{\footnotesize\color{inkfaint}\sffamily Star.org}
\fancyhead[R]{\footnotesize\color{inkfaint}\sffamily\WHTITLE}
\fancyfoot[C]{\footnotesize\color{inkfaint}\thepage}

% ---- 标题样式（不依赖 titlesec / sectsty）----
\makeatletter
\renewcommand{\section}{\@startsection{section}{1}{0pt}%
  {22pt plus 4pt minus 2pt}{10pt plus 2pt minus 2pt}%
  {\normalfont\sffamily\Large\bfseries\color{accent}}}
\renewcommand{\subsection}{\@startsection{subsection}{2}{0pt}%
  {16pt plus 3pt minus 2pt}{7pt plus 2pt minus 2pt}%
  {\normalfont\sffamily\large\bfseries\color{ink}}}
\renewcommand{\subsubsection}{\@startsection{subsubsection}{3}{0pt}%
  {12pt plus 2pt minus 2pt}{6pt plus 1pt minus 1pt}%
  {\normalfont\sffamily\normalsize\bfseries\color{inksoft}}}
\makeatother

% 章标题前不加首行缩进，段落间距由 parskip 控制
\setlength{\parindent}{0pt}

% ---- 引用块（核心主张）：左侧金色竖线 + 浅底 ----
\newcommand{\claim}[1]{%
  \par\vspace{6pt}%
  \noindent\begin{minipage}{\linewidth}%
    \setlength{\fboxsep}{12pt}%
    \colorbox{surface}{%
      \begin{minipage}{\dimexpr\linewidth-2\fboxsep\relax}%
        \color{accent}\bfseries\sffamily #1%
      \end{minipage}}%
  \end{minipage}\par\vspace{8pt}%
}

% ---- 一句话结语块 ----
\newcommand{\takeaway}[1]{%
  \par\vspace{8pt}\noindent\begin{minipage}{\linewidth}%
    \setlength{\fboxsep}{14pt}%
    \colorbox{accent}{%
      \begin{minipage}{\dimexpr\linewidth-2\fboxsep\relax}%
        \color{white}#1%
      \end{minipage}}%
  \end{minipage}\par\vspace{10pt}%
}

% ---- 表格统一样式 ----
\newcommand{\wtable}[1]{%
  \par\vspace{4pt}\noindent\small #1\par\vspace{6pt}%
}

% 目录条目字号
\renewcommand{\contentsname}{\sffamily\Large\bfseries\color{accent}\WHTOCNAME}

% 英文日期：ctexart 下 \today 输出中文格式（2026 年 9 月 21 日），英文版单独用这个
\newcommand{\WHdateEN}{%
  \ifcase\month\or January\or February\or March\or April\or May\or June\or
  July\or August\or September\or October\or November\or December\fi
  \space\the\day,\space\the\year}

% hyperref 书签与元数据
\hypersetup{
  colorlinks=true,
  linkcolor=accent,
  urlcolor=accent,
  pdftitle={\WHTITLE},
  pdfauthor={Star.org},
  pdfsubject={\WHSUBJECT},
}
，再写正文。
%  字体使用仓库内置的 Noto Sans/Serif SC 子集（fonts/），
%  渲染不依赖运行环境是否安装中文字体（见 DECISIONS D10）。
% =====================================================================

\usepackage{fontspec}
\usepackage{xeCJK}

% ---- 字体：仓库内置的 Noto Sans/Serif SC 子集 ----
% 路径相对于编译工作目录（docs/whitepaper/）。字体随仓库内置，
% 渲染结果不再依赖运行环境是否安装中文字体（DECISIONS D10）。
% 子集覆盖 CJK 统一表意文字 U+4E00–U+9FFF 全区间（见 fonts/manifest.json）。
% 内置子集只有 Regular 一个字重，AutoFakeBold 用于合成粗体，
% 否则 \bfseries 会静默回退到 Regular、标题看不出层级。
\setCJKmainfont[Path=../../fonts/,AutoFakeBold=2.5]{NotoSerifSC-Regular.subset.otf}
\setCJKsansfont[Path=../../fonts/,AutoFakeBold=2.5]{NotoSansSC-Regular.subset.otf}
\setCJKmonofont[Path=../../fonts/]{NotoSansSC-Regular.subset.otf}
% 拉丁字母与数字沿用 TeX 自带的 Latin Modern（任何 TeX Live 环境都有），
% 不依赖系统是否安装 Helvetica 一类的西文字体。

\usepackage{geometry}
\usepackage{xcolor}
\usepackage{booktabs}
\usepackage{array}
\usepackage{tabularx}
\usepackage{longtable}
\usepackage{fancyhdr}
\usepackage{setspace}
\usepackage{parskip}
\usepackage{caption}
\usepackage{ragged2e}
\usepackage{microtype}
\usepackage[hidelinks]{hyperref}

% ---- 页面 ----
\geometry{a4paper, top=26mm, bottom=24mm, left=24mm, right=24mm}

% ---- 配色（与站点一致的深蓝 + 金）----
\definecolor{ink}{HTML}{111827}
\definecolor{inksoft}{HTML}{4B5563}
\definecolor{inkfaint}{HTML}{6B7280}
\definecolor{accent}{HTML}{1E3A8A}
\definecolor{gold}{HTML}{B08430}
\definecolor{line}{HTML}{D8DEE9}
\definecolor{surface}{HTML}{F6F7FB}

% ---- 页眉页脚 ----
\pagestyle{fancy}
\fancyhf{}
\renewcommand{\headrulewidth}{0.4pt}
\renewcommand{\footrulewidth}{0pt}
\fancyhead[L]{\footnotesize\color{inkfaint}\sffamily Star.org}
\fancyhead[R]{\footnotesize\color{inkfaint}\sffamily\WHTITLE}
\fancyfoot[C]{\footnotesize\color{inkfaint}\thepage}

% ---- 标题样式（不依赖 titlesec / sectsty）----
\makeatletter
\renewcommand{\section}{\@startsection{section}{1}{0pt}%
  {22pt plus 4pt minus 2pt}{10pt plus 2pt minus 2pt}%
  {\normalfont\sffamily\Large\bfseries\color{accent}}}
\renewcommand{\subsection}{\@startsection{subsection}{2}{0pt}%
  {16pt plus 3pt minus 2pt}{7pt plus 2pt minus 2pt}%
  {\normalfont\sffamily\large\bfseries\color{ink}}}
\renewcommand{\subsubsection}{\@startsection{subsubsection}{3}{0pt}%
  {12pt plus 2pt minus 2pt}{6pt plus 1pt minus 1pt}%
  {\normalfont\sffamily\normalsize\bfseries\color{inksoft}}}
\makeatother

% 章标题前不加首行缩进，段落间距由 parskip 控制
\setlength{\parindent}{0pt}

% ---- 引用块（核心主张）：左侧金色竖线 + 浅底 ----
\newcommand{\claim}[1]{%
  \par\vspace{6pt}%
  \noindent\begin{minipage}{\linewidth}%
    \setlength{\fboxsep}{12pt}%
    \colorbox{surface}{%
      \begin{minipage}{\dimexpr\linewidth-2\fboxsep\relax}%
        \color{accent}\bfseries\sffamily #1%
      \end{minipage}}%
  \end{minipage}\par\vspace{8pt}%
}

% ---- 一句话结语块 ----
\newcommand{\takeaway}[1]{%
  \par\vspace{8pt}\noindent\begin{minipage}{\linewidth}%
    \setlength{\fboxsep}{14pt}%
    \colorbox{accent}{%
      \begin{minipage}{\dimexpr\linewidth-2\fboxsep\relax}%
        \color{white}#1%
      \end{minipage}}%
  \end{minipage}\par\vspace{10pt}%
}

% ---- 表格统一样式 ----
\newcommand{\wtable}[1]{%
  \par\vspace{4pt}\noindent\small #1\par\vspace{6pt}%
}

% 目录条目字号
\renewcommand{\contentsname}{\sffamily\Large\bfseries\color{accent}\WHTOCNAME}

% 英文日期：ctexart 下 \today 输出中文格式（2026 年 9 月 21 日），英文版单独用这个
\newcommand{\WHdateEN}{%
  \ifcase\month\or January\or February\or March\or April\or May\or June\or
  July\or August\or September\or October\or November\or December\fi
  \space\the\day,\space\the\year}

% hyperref 书签与元数据
\hypersetup{
  colorlinks=true,
  linkcolor=accent,
  urlcolor=accent,
  pdftitle={\WHTITLE},
  pdfauthor={Star.org},
  pdfsubject={\WHSUBJECT},
}
